\section{Alternatives and Justification}
\label{sec:alternatives}

This section first compares the transformation languages the course offers
(Section~\ref{sec:language}). It then compares the ways the transformation
could be structured, and gives the chosen solution in general terms
(Section~\ref{sec:structure}). Section~\ref{sec:design} describes the solution
in detail.

\subsection{The Transformation Language}
\label{sec:language}

The course offers ATL and QVT-Operational (QVTo). A third option is a
transformation written in Java, by hand or with a small rule library such as
SiTra. Table~\ref{tab:languages} compares the three against what this
transformation has to do.

\begin{table}[!htbp]
\centering
\small
\begin{tabular}{>{\raggedright\arraybackslash}p{0.27\linewidth}>{\raggedright\arraybackslash}p{0.20\linewidth}>{\raggedright\arraybackslash}p{0.20\linewidth}>{\raggedright\arraybackslash}p{0.20\linewidth}}
\hline
\textbf{Requirement} & \textbf{ATL} & \textbf{QVTo} & \textbf{Java} \\
\hline
Derive values from the whole cluster, not from one source element &
  helpers outside the rules & queries in the language & methods \\
Refer to a target object created by another rule &
  implicit resolution of matched rules &
  explicit \texttt{resolveone} through the trace & a trace kept by hand \\
Set a value only after the rest of the element exists &
  imperative \texttt{do} section & \texttt{end} section of a mapping & any order \\
Stop with a message on an underivable value &
  no assertion statement & \texttt{assert fatal} & exception \\
Call code for an operation the language lacks &
  native library only & Java black-box library & not needed \\
Graded as a transformation language & yes & yes & no \\
\hline
\end{tabular}
\caption{Transformation Languages Against the Requirements of the Resolution}
\label{tab:languages}
\end{table}

Most values in a resolved model do not come from one source element. The
eligible machines of a process depend on the node contract. A route's
precedence depends on the other routes of the same exposure. The ingress
rules of a process depend on every process in the cluster that depends on it.
ATL is strongest when one source element matches one target element, as in its
matched rules. Here, most of the work would end up in helpers and called rules,
which are imperative ATL without the matching. QVTo is imperative from the
start, uses OCL for its expressions, and makes the trace an explicit operation.
The transformation uses the trace in seven places, to point a gate member, a
route or a generated file at the process the transformation created for it.

A Java transformation would be the easiest to debug. It would also keep the
trace by hand, and it is not a transformation language, which the course
grades. QVTo was therefore chosen. The one operation QVTo cannot express,
computing a SHA-256 digest, is a Java black-box library
(Section~\ref{sec:blackbox}).

\subsection{The Structure of the Transformation}
\label{sec:structure}

Table~\ref{tab:structure} lists the structural alternatives that were
considered, and the choice made for each.

\begin{table}[!htbp]
\centering
\small
\begin{tabular}{>{\raggedright\arraybackslash}p{0.22\linewidth}>{\raggedright\arraybackslash}p{0.33\linewidth}>{\raggedright\arraybackslash}p{0.37\linewidth}}
\hline
\textbf{Question} & \textbf{Alternatives} & \textbf{Choice and reason} \\
\hline
One transformation or two &
  One transformation from the authored model; or a lowering first, then the resolution &
  Two. The constraints of Task~1 are checked on the lowered model, and the
  resolution reads one level of values instead of three. \\
How much to resolve per run &
  The whole cluster in one resolved model; or one project per run, with the
  other projects as input &
  One project per run. It matches the production implementation, which writes
  one resolved document per project, and keeps the model the templates read
  small. \\
Where the extra input documents come from &
  A Java YAML reader; hand-written XMI models; a metamodel with Xtext grammars &
  A metamodel with grammars. The same YAML files are read by both
  implementations, so their results can be compared. \\
How the inputs reach the transformation &
  One model extent per document; or one extent per metamodel &
  One extent for the intent (projects and platform) and one for the pinned
  inputs. The transformation selects documents by type. \\
How to compute digests &
  In OCL; in the Java pipeline before the run; in a black-box library &
  A black-box library. OCL cannot produce bytes, and a digest computed before
  the run could not cover values the transformation derives. \\
What to do with an underivable value &
  Write a default; write a diagnostic model; stop with a message &
  Stop. A default would produce files that deploy the wrong thing. \\
\hline
\end{tabular}
\caption{Structural Alternatives and the Choices Made}
\label{tab:structure}
\end{table}

\subsubsection{A Lowering Before the Resolution}
\label{sec:twotransformations}

An author may declare a secret, a dependency or a memory limit once for a
whole project, once for an application, or on a single process. The lowering
merges these levels onto each process, so that the \emph{effective} process
contains everything that applies to it. Task~1 checks its constraints on this
lowered model. If the resolution read the authored model instead, every rule
that reads a process would have to merge three levels again, and could merge
them differently from the constraints. Keeping two transformations means each
rule of the resolution reads one place.

\subsubsection{One Project per Run}
\label{sec:perproject}

The transformation runs once per project. The other projects are inputs: the
resolution of \texttt{auth} reads the \texttt{data} project to find the address
and port of the database it depends on. The alternative, one resolved model
for the whole cluster, would also work, because QVTo can create several
deployments in one run. It was rejected because the production implementation
produces one resolved document per project and the oracles are per project,
and because the templates of Task~3 generate one project's files at a time.

\subsubsection{Grammars for the Pinned Inputs}
\label{sec:pinnedalternatives}

The node contract, the images lock and the snapshot are YAML files in the same
subset as the project documents. Three ways of reading them were compared. A
Java YAML reader would be short, but it would add a parser the course does not
teach, and it would read the files differently from the Xtext grammars that
read the project documents. Hand-written XMI models of the three documents
would be the quickest, but then the transformation would never read the files
the production implementation reads, and the comparison between the two
implementations would not cover them. The chosen solution adds one grammar per
document. Each grammar inherits the terminals and indentation handling of the
project grammar, so it adds only its own rules. The migration proof is read the
same way. An Asset's file is the one input that is not parsed: its text belongs
to the application, so the pipeline reads it as a string.

\subsubsection{Stopping Instead of Defaulting}
\label{sec:stopping}

The specification states every derivation as total: each value either derives
from declared input or the input is rejected. Some rejections need several
documents, for example an image that the lock does not contain, and the
constraints of Task~1 do not check those yet. The transformation therefore
contains fatal assertions at each point where it would otherwise have to
invent a value. Each assertion names what is missing. The same mechanism let
the transformation be extended one example at a time: while an example was not
yet supported, the transformation stopped on it with the name of the ticket
that would add it.
